\documentclass[12pt,a4paper]{article}
\usepackage[T1]{fontenc}
\usepackage{lmodern}
\usepackage{textcomp}
\usepackage[a4paper,margin=1in,headheight=15pt]{geometry}
\emergencystretch=3em
\usepackage{graphicx}
\usepackage{longtable,booktabs,array,calc}
\usepackage{float}
\usepackage{caption}
\usepackage{enumitem}
\usepackage[expansion=false]{microtype}
\usepackage{fancyhdr}
\usepackage{hyperref}
\hypersetup{colorlinks=true,linkcolor=black,urlcolor=blue,citecolor=black}
\pagestyle{fancy}
\fancyhf{}
\fancyfoot[C]{\thepage}
\renewcommand{\headrulewidth}{0pt}
\setlength{\parindent}{0pt}
\setlength{\parskip}{5pt plus 1pt minus 1pt}
\setcounter{secnumdepth}{0}
\setcounter{tocdepth}{2}
\renewcommand{\arraystretch}{1.18}
\setlength{\LTpre}{6pt}
\setlength{\LTpost}{6pt}
\providecommand{\tightlist}{\setlength{\itemsep}{2pt}\setlength{\parskip}{0pt}}
\providecommand{\real}[1]{#1}
\begin{document}
\sloppy
\thispagestyle{fancy}
\begin{center}
\vspace*{0.05cm}
{\Large\bfseries INDIAN INSTITUTE OF INFORMATION TECHNOLOGY\par}
\vspace{0.05cm}
{\Large\bfseries VADODARA\par}
\vspace{0.25cm}
\includegraphics[width=1.1in,height=1.1in]{media/media/image1.png}\par
\vspace{0.25cm}
{\Large\bfseries CS545\par}
\vspace{0.06cm}
{\Large\bfseries SOFTWARE ENGINEERING\par}
\vspace{0.06cm}
{\LARGE\bfseries PROJECT REPORT\par}
\vspace{0.2cm}
{\large Master of Computer Applications (MCA)\par}
{\large Second Year\par}
\vspace{0.2cm}
{\large\bfseries Project Title:\par}
\vspace{0.12cm}
{\large\bfseries Wereverse: Unified Retail Management System for Shop Billing, Khata, WhatsApp Bills,\par}
{\large\bfseries Tally Integration, Payroll and Saree Fall-Pico\par}
\vspace{0.15cm}
{\large\bfseries Group No. 36\par}
\vspace{0.15cm}
{\large\bfseries Submitted By:\par}
\vspace{0.08cm}
\renewcommand{\arraystretch}{1.05}
\begin{tabular}{c l c}
\textbf{S. No.} & \textbf{Name} & \textbf{Enrollment ID}\\
1 & Ankit Kumar Sharma & 20252651007\\
2 & Piyush Lalchand Menghani & 20252651039\\
3 & Shubham Gupta & 20252651021\\
4 & Vijay Singh & 20252651064\\
5 & Vivek Kumar Kosta & 20252651065\\
\end{tabular}
\vspace{0.15cm}

{\large\bfseries Submitted To:\par}
\vspace{0.04cm}
Faculty Name: Rishabh Gupta\par
Department: MCA\par
\vspace{0.12cm}
{\large\bfseries Indian Institute of Information Technology\par}
{\large\bfseries Vadodara Autumn 2026-27\par}
\end{center}
\clearpage
\section*{CERTIFICATE}\label{certificate}

This is to certify that the following students of the \textbf{Master of
Computer Applications (MCA)} programme at the \textbf{Indian Institute
of Information Technology Vadodara} have successfully completed the
laboratory work for \textbf{CS545 -- Software Engineering}, covering
Experiments 1 to 9, during the academic year \textbf{2026--2027}.

The laboratory work has been carried out as part of the course
requirements through the team project titled \textbf{``Wereverse''.}

Wereverse is a unified retail management system for a fashion shop. It
combines shop billing that works offline, one shared stock for shop and
website, khata (credit) with WhatsApp bills and a Pay Now link, two-way
Tally integration, employee payroll, and saree fall-pico job tracking.

The details of the team members are as follows:

\begin{longtable}[]{@{}
  >{\raggedright\arraybackslash}p{(\columnwidth - 4\tabcolsep) * \real{0.1217}}
  >{\raggedright\arraybackslash}p{(\columnwidth - 4\tabcolsep) * \real{0.4956}}
  >{\raggedright\arraybackslash}p{(\columnwidth - 4\tabcolsep) * \real{0.3827}}@{}}
\toprule\noalign{}
\begin{minipage}[b]{\linewidth}\raggedright
\textbf{S. No.}
\end{minipage} & \begin{minipage}[b]{\linewidth}\raggedright
\textbf{Name of Student}
\end{minipage} & \begin{minipage}[b]{\linewidth}\raggedright
\textbf{Enrollment Number}
\end{minipage} \\
\midrule\noalign{}
\endhead
\bottomrule\noalign{}
\endlastfoot
1 & Ankit Kumar Sharma & 20252651007 \\
2 & Piyush Lalchand Menghani & 20252651039 \\
3 & Shubham Gupta & 20252651021 \\
4 & Vijay Singh & 20252651064 \\
5 & Vivek Kumar Kosta & 20252651065 \\
\end{longtable}

This certificate is issued in recognition of the completion and
submission of the prescribed laboratory work in accordance with the
requirements of the course.

\begin{quote}
Date: 29-10-2026

Faculty Signature: \_\_\_\_\_\_\_\_\_\_\_\_\_\_\_\_\_\_\_\_

Faculty Name: Rishabh Gupta
\end{quote}

\clearpage
\section*{DECLARATION}\label{declaration}

\begin{quote}
We hereby declare that the project entitled ``Wereverse'', submitted as
part of the coursework for CS545 -- Software Engineering at the Indian
Institute of Information Technology Vadodara, is our original work
carried out under the guidance of Rishabh Gupta.

We declare that this report has been prepared by our team for academic
purposes. The information presented in this report is accurate to the
best of our knowledge, and all sources used have been appropriately
acknowledged. We understand that any work or material taken from other
sources must be properly cited.

\textbf{Project Title:} Wereverse

\textbf{Course Code:} CS545 -- Software Engineering

\textbf{Group Number:} 36

\textbf{Team Members:}
\end{quote}

\begin{longtable}[]{@{}
  >{\raggedright\arraybackslash}p{(\columnwidth - 4\tabcolsep) * \real{0.1217}}
  >{\raggedright\arraybackslash}p{(\columnwidth - 4\tabcolsep) * \real{0.4956}}
  >{\raggedright\arraybackslash}p{(\columnwidth - 4\tabcolsep) * \real{0.3827}}@{}}
\toprule\noalign{}
\begin{minipage}[b]{\linewidth}\raggedright
\textbf{S. No.}
\end{minipage} & \begin{minipage}[b]{\linewidth}\raggedright
\textbf{Name of Student}
\end{minipage} & \begin{minipage}[b]{\linewidth}\raggedright
\textbf{Enrollment Number}
\end{minipage} \\
\midrule\noalign{}
\endhead
\bottomrule\noalign{}
\endlastfoot
1 & Ankit Kumar Sharma & 20252651007 \\
2 & Piyush Lalchand Menghani & 20252651039 \\
3 & Shubham Gupta & 20252651021 \\
4 & Vijay Singh & 20252651064 \\
5 & Vivek Kumar Kosta & 20252651065 \\
\end{longtable}

\begin{quote}
\textbf{Place:} IIIT Vadodara

\textbf{Date:} 29-10-2026

\textbf{Signatures of Team Members:}

1. Ankit Kumar Sharma: \_\_\_\_\_\_\_\_\_\_\_\_\_\_\_\_\_\_\_\_

2. Piyush Lalchand Menghani: \_\_\_\_\_\_\_\_\_\_\_\_\_\_\_\_\_\_\_\_

3. Shubham Gupta: \_\_\_\_\_\_\_\_\_\_\_\_\_\_\_\_\_\_\_\_

4. Vijay Singh: \_\_\_\_\_\_\_\_\_\_\_\_\_\_\_\_\_\_\_\_

5. Vivek Kumar Kosta: \_\_\_\_\_\_\_\_\_\_\_\_\_\_\_\_\_\_\_\_
\end{quote}

\clearpage
\section*{ACKNOWLEDGEMENT}\label{acknowledgement}

\begin{quote}
We would like to express our sincere gratitude to Rishabh Gupta, our
faculty guide for the course CS545 -- Software Engineering, for his
guidance and support throughout the development of our project,
``Wereverse.''

We are thankful to the Indian Institute of Information Technology
Vadodara for providing us with the opportunity to work on this project
and apply our software engineering concepts in a practical setting.

We would also like to thank our faculty members and classmates for their
valuable suggestions and encouragement. Their support has helped us
improve our understanding of software development, system design, and
project management.

Finally, we would like to thank all our team members for their
cooperation, teamwork, and contributions toward the planning and
development of Wereverse.

\textbf{Team Members:}

1. Ankit Kumar Sharma

2. Piyush Lalchand Menghani

3. Shubham Gupta

4. Vijay Singh

5. Vivek Kumar Kosta

\textbf{Group No.:} 36
\end{quote}

\clearpage
\section*{TABLE OF CONTENTS}\label{table-of-contents}
\vspace{0.5cm}
\begin{center}
\renewcommand{\arraystretch}{1.55}
\begin{tabular}{|c|p{10.2cm}|c|}
\hline
\textbf{Experiment No.} & \centering\arraybackslash\textbf{Title} & \textbf{Page No.}\\
\hline
1 & Problem Statement and Software Requirements Specification (SRS) & \pageref{lab-experiment-1}\\
\hline
2 & Software Process Model, Project Planning and Risk Management & \pageref{lab-experiment-2}\\
\hline
3 & Use Case Diagram and Use Case Descriptions & \pageref{lab-experiment-3}\\
\hline
4 & Class Diagram & \pageref{lab-experiment-4}\\
\hline
5 & Sequence Diagrams & \pageref{lab-experiment-5}\\
\hline
6 & Activity Diagram and State Chart Diagram & \pageref{lab-experiment-6}\\
\hline
7 & ER Diagram, Deployment Diagram and Test Cases & \pageref{lab-experiment-7}\\
\hline
8 & Modular Design of a Software System & \pageref{lab-experiment-8}\\
\hline
9 & Functional Test Case Design & \pageref{lab-experiment-9}\\
\hline
 & Overall Conclusion & \pageref{overall-conclusion}\\
\hline
 & References & \pageref{references}\\
\hline
\end{tabular}
\end{center}
\clearpage

\section{LAB EXPERIMENT 1}\label{lab-experiment-1}

\subsection{Problem Statement and Software Requirements Specification
(SRS)}\label{problem-statement-and-software-requirements-specification-srs}

\begin{quote}
\textbf{Aim}

To study the problem domain of Wereverse and prepare its Software
Requirements Specification: problem statement, scope, users, functional
requirements and non-functional requirements.
\end{quote}

\subsection{Problem Statement}\label{problem-statement}

\begin{quote}
The client runs one fashion shop (with saree sales) and wants to sell
through a website as well. Today the counter bills in TallyPrime, credit
sales (\emph{khata}) are tracked by hand, customers are not told their
dues quickly, staff salary and advances are calculated manually, and
saree fall-pico work given to outside workers is tracked in a register.
Shop and website stock are not linked, so items can be oversold and
counted twice.

\textbf{Wereverse} is one system that runs the whole business: shop
billing that works offline, one shared stock for shop and website, khata
with instant WhatsApp bills and a Pay Now link, two-way Tally
integration so the accountant never enters data twice, employee payroll,
and fall-pico job and wage tracking. AI features (recommendations,
chatbot, demand forecast) are added in later phases.
\end{quote}

\subsection{Objectives}\label{objectives}

\begin{itemize}
\item
  One stock for shop and website, so nothing is oversold and nobody
  counts twice.
\item
  Fast shop billing that keeps working when the internet is down.
\item
  Khata customers get an instant WhatsApp bill and can pay from home;
  the owner always knows who owes what.
\item
  Tally always matches Wereverse, without double entry by the
  accountant.
\item
  Clear staff salary, advance and fall-pico wage records so disputes are
  rare.
\item
  AI to increase sales and improve buying decisions (Phase 2 and 3).
\end{itemize}

\subsection{Scope}\label{scope}

\begin{longtable}[]{@{}
  >{\raggedright\arraybackslash}p{(\columnwidth - 2\tabcolsep) * \real{0.6000}}
  >{\raggedright\arraybackslash}p{(\columnwidth - 2\tabcolsep) * \real{0.4000}}@{}}
\toprule\noalign{}
\begin{minipage}[b]{\linewidth}\raggedright
\textbf{In scope (Phase 1)}
\end{minipage} & \begin{minipage}[b]{\linewidth}\raggedright
\textbf{Out of scope (for now)}
\end{minipage} \\
\midrule\noalign{}
\endhead
\bottomrule\noalign{}
\endlastfoot
Shared catalog and stock with barcodes; POS billing with offline mode;
khata with WhatsApp bills and Pay Now; two-way Tally integration;
employees and payroll; fall-pico; basic purchasing; basic website
(browse, cart, checkout, payments, shipping, returns); roles, audit log,
owner dashboard and reports. & Native mobile apps, international sales,
wholesale portal, AR try-on, Amazon/Flipkart sync, multiple shops, and a
SaaS platform for other brands. \\
\end{longtable}

\subsection{Users of the System}\label{users-of-the-system}

\begin{longtable}[]{@{}
  >{\raggedright\arraybackslash}p{(\columnwidth - 2\tabcolsep) * \real{0.3200}}
  >{\raggedright\arraybackslash}p{(\columnwidth - 2\tabcolsep) * \real{0.6800}}@{}}
\toprule\noalign{}
\begin{minipage}[b]{\linewidth}\raggedright
\textbf{User}
\end{minipage} & \begin{minipage}[b]{\linewidth}\raggedright
\textbf{Main needs}
\end{minipage} \\
\midrule\noalign{}
\endhead
\bottomrule\noalign{}
\endlastfoot
Owner & Full view of sales, cash, khata, stock and salary; approvals;
buying decisions \\
Manager & Daily operations, discounts, attendance, fall-pico
assignment \\
Cashier / sales staff & Quick billing, khata sales, returns \\
Accountant & Tally sync, reconciliation, GST, salary sheets \\
Online order staff & Pick, pack and ship website orders \\
Fall-pico supervisor and workers & Job list, quality check, earnings \\
Employee & Own attendance, salary, advance and balance \\
Khata customer & Own bills and balance; pay from home \\
Online customer & Browse, buy, track and return \\
\end{longtable}

\subsection{Functional Requirements}\label{functional-requirements}

\begin{quote}
Requirements are grouped by SRS module code. Each module is delivered in
Phase 1 unless stated.
\end{quote}

\begin{longtable}[]{@{}
  >{\raggedright\arraybackslash}p{(\columnwidth - 2\tabcolsep) * \real{0.2800}}
  >{\raggedright\arraybackslash}p{(\columnwidth - 2\tabcolsep) * \real{0.7200}}@{}}
\toprule\noalign{}
\begin{minipage}[b]{\linewidth}\raggedright
\textbf{Module (SRS ID)}
\end{minipage} & \begin{minipage}[b]{\linewidth}\raggedright
\textbf{Functional requirements}
\end{minipage} \\
\midrule\noalign{}
\endhead
\bottomrule\noalign{}
\endlastfoot
Catalog and stock (FR-CAT) & Maintain products and variants (price, MRP,
GST rate, HSN, fall-pico flag); print barcode labels; keep one stock
ledger where every change is a stock movement; reserve stock for unpaid
online orders and release after 30 minutes; stock count with reasoned
adjustments. \\
Shop billing (FR-POS) & Scan-first billing; split payment (cash, UPI,
card, khata); discounts within role limits, above limit needs manager
PIN; returns and exchanges; cash register open/close; works offline with
a local queue and idempotent sync (FR-POS-08). \\
Khata (FR-KHT) & Open khata accounts with credit limit and credit days;
append-only ledger (sale, payment, return, adjustment, write-off,
reversal); credit-limit check with manager override; oldest-bill-first
allocation; ageing report; reminder rules. \\
WhatsApp (FR-WAP) & Send bill, receipt, reminder and fall-pico-ready
messages using approved templates; AI writes wording around placeholders
only; validator rejects any wrong number; SMS fallback; opt-in, opt-out
and quiet hours. \\
Pay from home (FR-PRT, FR-PAY) & Signed Pay Now link valid 7 days
showing balance and bills; part or full payment through Razorpay;
payment confirmed only by verified webhook; overpayment becomes advance
credit. \\
Tally integration (FR-TLY) & Mode A: import bills made in TallyPrime;
Mode B: push POS sales to Tally; always push online orders and receipts;
mapping of items and ledgers; Sync Issues screen; nightly
reconciliation; unique voucher reference so retries never duplicate. \\
Employees (FR-EMP) & Attendance (manager entry + PIN); monthly and daily
salary; incentive, overtime, advances and recovery; salary run draft,
owner approval and lock; employee sees own statement only. \\
Fall-pico (FR-FPC) & Create a job for each eligible saree above the
price limit (default Rs. 1,000); print tag; assign workers; QC pass/fail
with rework; per-piece wage on QC pass; online order cannot be packed
until QC passed or skipped. \\
Purchasing (FR-PUR-01..03) & Supplier bills, stock-in and label
printing. \\
Website (FR-STO, FR-CHK, FR-ORD) & Browse, search, cart, checkout, order
tracking, shipping and returns; no khata option for online customers. \\
Admin and reports (FR-AUT, FR-ADM, FR-RPT, FR-NOT) & Role-based access,
complete audit log, settings, owner dashboard, sales and GST reports,
notifications. \\
AI (FR-AI, Phase 2/3) & Recommendations, chatbot, smart search, size
help, description writer, poster maker; later visual search, demand
forecast and what-to-buy. \\
\end{longtable}

\subsection{Non-Functional
Requirements}\label{non-functional-requirements}

\begin{longtable}[]{@{}
  >{\raggedright\arraybackslash}p{(\columnwidth - 2\tabcolsep) * \real{0.2400}}
  >{\raggedright\arraybackslash}p{(\columnwidth - 2\tabcolsep) * \real{0.7600}}@{}}
\toprule\noalign{}
\begin{minipage}[b]{\linewidth}\raggedright
\textbf{Category}
\end{minipage} & \begin{minipage}[b]{\linewidth}\raggedright
\textbf{Requirement}
\end{minipage} \\
\midrule\noalign{}
\endhead
\bottomrule\noalign{}
\endlastfoot
Performance & Item appears within 1 second of scan; 3-item cash bill
under 30 seconds; API p95 under 400 ms; website LCP under 2.5 s on
4G. \\
Availability / offline & POS keeps billing with no internet; queue
clears within 1 minute of internet returning; AI failure never stops
billing. \\
Data integrity & Money stored as whole paise; ledgers and stock tables
append-only; daily integrity job recomputes balances (NFR-DATA-01); zero
duplicate or lost Tally vouchers. \\
Security & HTTPS, hashed passwords, 2FA for owner; role checks on every
API; webhooks verified by signature; Tally interface never exposed to
the internet. \\
Privacy & Salary and khata data encrypted and access-logged; no personal
data sent to external AI models; WhatsApp opt-in recorded. \\
Usability & POS text at least 16 px and main buttons at least 56 px;
simple Hindi-friendly pages for workers and customers; mobile-first
website. \\
Maintainability & Modular monolith (FastAPI), typed code, migrations
only through Alembic, 2 reviewers for money and Tally code. \\
\end{longtable}

\subsection{Constraints and
Assumptions}\label{constraints-and-assumptions}

\begin{itemize}
\item
  The shop has one counter and bills in TallyPrime today; Wereverse
  starts in Mode A and moves to its own POS (Mode B) after a clean
  parallel run.
\item
  TallyPrime runs on a 24-hour cloud Windows server in India, with the
  sync agent on the same server.
\item
  Technology stack: FastAPI, React + Vite + TypeScript, PostgreSQL,
  Celery + Redis, hosted on AWS Mumbai.
\item
  Items marked {[}Assumption{]} or {[}TBD{]} in the source documentation
  (hardware models, WhatsApp provider, hosting licence) are still to be
  confirmed by the client.
\end{itemize}

\subsection{Conclusion}\label{conclusion}

\begin{quote}
The SRS converts the shop\textquotesingle s daily problems into clear,
testable requirements grouped by module. Separating functional from
non-functional requirements made the hard limits (offline billing, exact
money, no duplicate vouchers) visible early, and these drive every later
design decision.
\end{quote}

\clearpage
\section{LAB EXPERIMENT 2}\label{lab-experiment-2}

\subsection{Software Process Model, Project Planning and Risk
Management}\label{software-process-model-project-planning-and-risk-management}

\begin{quote}
\textbf{Aim}

To select a suitable software process model for Wereverse, and to
prepare the project plan: team roles, schedule, milestones and risk
management.
\end{quote}

\subsection{Choice of Process Model: Incremental Model with Agile (Scrum)
Sprints}\label{choice-of-process-model-incremental-model-with-agile-scrum-sprints}

\begin{quote}
Wereverse is built as a series of increments, each delivered through
two-week Scrum sprints. Each increment is used in the real shop for
about a week before the next one starts. This model was selected
because:
\end{quote}

\begin{itemize}
\item
  \textbf{The shop must never stop billing.} The counter keeps billing
  in TallyPrime (Mode A) while new modules go live one by one; the POS
  replaces Tally billing only after a clean parallel run.
\item
  \textbf{Requirements are mostly known but some decisions are open}
  (hardware, WhatsApp provider, fall-pico rules). Incremental delivery
  lets these be settled without blocking earlier modules.
\item
  \textbf{Early value:} khata with WhatsApp bills and Pay Now reaches
  the shop by Week 10, long before the website.
\item
  \textbf{Risk control:} the riskiest part (Tally sync) is built and
  tested first, against a test company.
\item
  A pure Waterfall model would deliver nothing until the end; a pure
  prototype model does not suit money and GST logic that must be exact.
\end{itemize}

\subsubsection{Rollout order
(increments)}\label{rollout-order-increments}

\begin{itemize}
\item
  Increment 1: data import, shared stock and Tally sync while the
  counter bills in TallyPrime (Mode A).
\item
  Increment 2: khata, WhatsApp bills and Pay Now.
\item
  Increment 3: POS billing run side by side with Tally, then cutover to
  Mode B.
\item
  Increment 4: employees, fall-pico and purchasing.
\item
  Increment 5: the website. AI features follow as Phase 2 and Phase 3.
\end{itemize}

\subsection{Team Roles}\label{team-roles}

\begin{longtable}[]{@{}
  >{\raggedright\arraybackslash}p{(\columnwidth - 4\tabcolsep) * \real{0.2700}}
  >{\raggedright\arraybackslash}p{(\columnwidth - 4\tabcolsep) * \real{0.4800}}
  >{\raggedright\arraybackslash}p{(\columnwidth - 4\tabcolsep) * \real{0.2500}}@{}}
\toprule\noalign{}
\begin{minipage}[b]{\linewidth}\raggedright
\textbf{Role}
\end{minipage} & \begin{minipage}[b]{\linewidth}\raggedright
\textbf{Owns}
\end{minipage} & \begin{minipage}[b]{\linewidth}\raggedright
\textbf{Also helps with}
\end{minipage} \\
\midrule\noalign{}
\endhead
\bottomrule\noalign{}
\endlastfoot
Frontend and POS developer & POS app (offline), back-office UI, website,
simple pages & Design QA \\
Backend developer (lead) & Core API, stock, billing, khata ledger,
payroll, fall-pico, database, hosting & CI/CD \\
Integration developer & Tally sync agent, WhatsApp, Razorpay, courier,
SMS and email & Reconciliation and reports \\
AI engineer & AI service, WhatsApp wording and validator,
recommendations, chatbot, forecasts & Data pipelines, search \\
PM, QA and design lead & SRS, planning, Figma, test cases, UAT, client
and accountant communication & Training and handover \\
\end{longtable}

\subsection{Project Schedule (Phase 0 and Phase 1, 20
weeks)}\label{project-schedule-phase-0-and-phase-1-20-weeks}

\includegraphics[width=6.235in,height=2.32906in]{media/media/image2.jpeg}

\begin{quote}
\emph{Figure 2.1: Gantt chart of sprints S0 to S9 with milestones M1 to
M6}
\end{quote}

\begin{longtable}[]{@{}
  >{\raggedright\arraybackslash}p{(\columnwidth - 4\tabcolsep) * \real{0.1400}}
  >{\raggedright\arraybackslash}p{(\columnwidth - 4\tabcolsep) * \real{0.1000}}
  >{\raggedright\arraybackslash}p{(\columnwidth - 4\tabcolsep) * \real{0.7600}}@{}}
\toprule\noalign{}
\begin{minipage}[b]{\linewidth}\raggedright
\textbf{Milestone}
\end{minipage} & \begin{minipage}[b]{\linewidth}\raggedright
\textbf{Week}
\end{minipage} & \begin{minipage}[b]{\linewidth}\raggedright
\textbf{Deliverable}
\end{minipage} \\
\midrule\noalign{}
\endhead
\bottomrule\noalign{}
\endlastfoot
M1 & 2 & SRS signed; designs and architecture approved; Tally cloud
server access and version confirmed; Tally data sample checked \\
M2 & 6 & Mode A on staging: Tally import (items, parties, vouchers),
shared stock, Tally sync, Sync Issues screen \\
M3 & 10 & Khata, WhatsApp bills, Pay Now and Tally receipts live in the
shop \\
M4 & 16 & POS billing (offline), employees, purchasing and fall-pico
working; POS parallel run in progress \\
M5 & 18 & Website checkout, orders and shipping working; UAT done \\
M6 & 20 & Parallel run clean, new bill series approved, POS cutover and
go-live \\
\end{longtable}

\subsection{Risk Management}\label{risk-management}

\begin{longtable}[]{@{}
  >{\raggedright\arraybackslash}p{(\columnwidth - 8\tabcolsep) * \real{0.0700}}
  >{\raggedright\arraybackslash}p{(\columnwidth - 8\tabcolsep) * \real{0.2700}}
  >{\raggedright\arraybackslash}p{(\columnwidth - 8\tabcolsep) * \real{0.1600}}
  >{\raggedright\arraybackslash}p{(\columnwidth - 8\tabcolsep) * \real{0.1500}}
  >{\raggedright\arraybackslash}p{(\columnwidth - 8\tabcolsep) * \real{0.3500}}@{}}
\toprule\noalign{}
\textbf{\#} & \textbf{Risk} & \textbf{Impact} & \textbf{Likelihood} & \textbf{Mitigation} \\
\midrule\noalign{}
\endhead
\bottomrule\noalign{}
\endlastfoot
R1 & Tally sync harder than expected (version, locked periods,
duplicates) & Wrong books, delay & High & Check in Sprint 0; test
company; outbox with unique references; daily reconciliation \\
R2 & WhatsApp template approval slow or rejected & Bills cannot be sent
& Medium & Submit in Week 1; SMS fallback; simple wording \\
R3 & Scope too large for 5 people & Late launch & High & Phase plan;
shop first, website last; AI after go-live; strict change control \\
R4 & Offline POS creates conflicts or duplicate bills & Wrong stock or
money & Medium & Idempotency keys, per-counter series, conflict alerts,
offline testing \\
R5 & Wrong opening balances at migration & Customer disputes & Medium &
Owner sign-off per balance; parallel run \\
R8 & AI writes a wrong amount & Trust and money & Low & Placeholders,
validator, template fallback, zero-error test \\
R13 & Counter billing in cloud Tally stops when internet fails & Billing
stops & Medium & Backup internet, paper bill book, move to offline
POS \\
R15 & Tally data quality (duplicate parties, missing phones) & Wrong or
missing messages & Medium & Sample check in Sprint 0; clean-up; Sync
Issues screen \\
R17 & New bill number series wrong under GST & Compliance & Low &
Design with accountant; test on staging; rehearse cutover \\
\end{longtable}

\subsection{Conclusion}\label{conclusion-1}

\begin{quote}
The incremental model with two-week sprints fits a live business that
cannot stop trading. The plan front-loads the riskiest integration
(Tally), delivers khata value early, and uses a parallel run as the gate
before the POS cutover. The risk register links each major risk to a
concrete mitigation.
\end{quote}

\clearpage
\section{LAB EXPERIMENT 3}\label{lab-experiment-3}

\subsection{Use Case Diagram and Use Case
Descriptions}\label{use-case-diagram-and-use-case-descriptions}

\begin{quote}
\textbf{Aim}

To identify the actors and use cases of Wereverse and draw the use case
diagram, with detailed descriptions for the key use cases.
\end{quote}

\subsection{Actors}\label{actors}

\begin{longtable}[]{@{}
  >{\raggedright\arraybackslash}p{(\columnwidth - 4\tabcolsep) * \real{0.2700}}
  >{\raggedright\arraybackslash}p{(\columnwidth - 4\tabcolsep) * \real{0.1800}}
  >{\raggedright\arraybackslash}p{(\columnwidth - 4\tabcolsep) * \real{0.5500}}@{}}
\toprule\noalign{}
\begin{minipage}[b]{\linewidth}\raggedright
\textbf{Actor}
\end{minipage} & \begin{minipage}[b]{\linewidth}\raggedright
\textbf{Type}
\end{minipage} & \begin{minipage}[b]{\linewidth}\raggedright
\textbf{Description}
\end{minipage} \\
\midrule\noalign{}
\endhead
\bottomrule\noalign{}
\endlastfoot
Cashier & Primary & Makes cash, split and khata bills, returns and day
close at the counter \\
Manager & Primary & Approves over-limit credit and discounts; marks
attendance \\
Owner & Primary & Views dashboard and reports; approves salary runs,
advances and write-offs \\
Accountant & Primary & Monitors Tally sync, fixes sync issues,
reconciles, prepares salary sheets \\
Employee & Primary & Views own attendance, salary and advance balance \\
Khata customer & Primary & Receives WhatsApp bills and pays dues through
the Pay Now link \\
Online customer & Primary & Browses, buys, tracks and returns orders on
the website \\
Fall-pico supervisor / worker & Primary & Assigns jobs and checks
quality / views own jobs and earnings \\
WhatsApp API, Razorpay, Tally agent & Secondary (systems) & External
systems for messages, payments and accounting sync \\
\end{longtable}

\begin{quote}
\includegraphics[width=5.7686in,height=6.81406in]{media/media/image3.jpeg}

\emph{Figure 3.1: Use case diagram of Wereverse}
\end{quote}

\subsection{Use Case Description: Make Khata (Credit)
Bill}\label{use-case-description-make-khata-credit-bill}

\begin{longtable}[]{@{}
  >{\raggedright\arraybackslash}p{(\columnwidth - 2\tabcolsep) * \real{0.2500}}
  >{\raggedright\arraybackslash}p{(\columnwidth - 2\tabcolsep) * \real{0.7500}}@{}}
\toprule\noalign{}
\begin{minipage}[b]{\linewidth}\raggedright
\textbf{Field}
\end{minipage} & \begin{minipage}[b]{\linewidth}\raggedright
\textbf{Description}
\end{minipage} \\
\midrule\noalign{}
\endhead
\bottomrule\noalign{}
\endlastfoot
Use case ID / Name & UC-02 / Make khata (credit) bill \\
Primary actor & Cashier \\
Secondary actors & Manager, WhatsApp API, Tally agent \\
Pre-conditions & Cashier is logged in; customer has an open khata
account; items have barcodes. \\
Main flow & 1. Cashier finds customer by phone. 2. Scans items. 3.
Chooses Khata as payment. 4. System shows balance and limit (include:
Check credit limit). 5. Cashier confirms. 6. System saves the bill,
deducts stock, adds a SALE ledger entry and queues WhatsApp and Tally
messages. 7. Bill is printed and the customer receives a WhatsApp bill
with a Pay Now link. \\
Alternate flows & 4a. Limit crossed: manager approves with PIN and
reason (extend: Approve over-limit), otherwise payment mode is changed.
6a. No internet: bill is saved in the local queue and synced later with
the same idempotency key. 7a. WhatsApp fails: SMS is sent and the entry
is flagged. \\
Post-conditions & Ledger balance increased; stock reduced; fall-pico
jobs created for eligible sarees; Tally sales voucher posted (Mode
B). \\
Related requirements & FR-KHT, FR-POS-08, FR-WAP, FR-TLY \\
\end{longtable}

\subsection{Use Case Description: Pay Dues via Pay Now
Link}\label{use-case-description-pay-dues-via-pay-now-link}

\begin{longtable}[]{@{}
  >{\raggedright\arraybackslash}p{(\columnwidth - 2\tabcolsep) * \real{0.2500}}
  >{\raggedright\arraybackslash}p{(\columnwidth - 2\tabcolsep) * \real{0.7500}}@{}}
\toprule\noalign{}
\begin{minipage}[b]{\linewidth}\raggedright
\textbf{Field}
\end{minipage} & \begin{minipage}[b]{\linewidth}\raggedright
\textbf{Description}
\end{minipage} \\
\midrule\noalign{}
\endhead
\bottomrule\noalign{}
\endlastfoot
Use case ID / Name & UC-08 / Pay dues via Pay Now link \\
Primary actor & Khata customer \\
Secondary actors & Razorpay, WhatsApp API, Tally agent \\
Pre-conditions & Customer received a valid signed link (valid 7
days). \\
Main flow & 1. Customer opens the link. 2. System shows balance and
recent bills. 3. Customer chooses an amount and pays by UPI. 4. Razorpay
sends a signed webhook. 5. System writes a PAYMENT ledger entry,
allocates it to the oldest bills, and queues a WhatsApp receipt and a
Tally receipt voucher. \\
Alternate flows & 1a. Link expired or belongs to someone else: access
denied. 4a. Duplicate webhook: ignored. 5a. Overpayment: stored as
advance credit. \\
Post-conditions & Balance reduced; receipt sent; Tally updated. \\
Related requirements & FR-PRT, FR-PAY, FR-KHT \\
\end{longtable}

\subsection{Conclusion}\label{conclusion-2}

\begin{quote}
The use case diagram shows who uses Wereverse and for what, across the
shop, the back office, the website and external systems. The include and
extend relationships capture the credit-limit rule, and the detailed
descriptions record the alternate flows (offline, over-limit, duplicate
webhook) that matter most for correctness.
\end{quote}

\clearpage
\section{LAB EXPERIMENT 4}\label{lab-experiment-4}

\subsection{Class Diagram}\label{class-diagram}

\begin{quote}
\textbf{Aim}

To identify the main classes of Wereverse with their attributes,
operations and relationships, and draw the class diagram.

\includegraphics[width=5.55764in,height=8.6075in]{media/media/image4.png}

\emph{Figure 4.1: Class diagram of Wereverse (core domain), rotated for
readability}
\end{quote}

\subsection{Class Responsibilities}\label{class-responsibilities}

\begin{longtable}[]{@{}
  >{\raggedright\arraybackslash}p{(\columnwidth - 2\tabcolsep) * \real{0.3000}}
  >{\raggedright\arraybackslash}p{(\columnwidth - 2\tabcolsep) * \real{0.7000}}@{}}
\toprule\noalign{}
\begin{minipage}[b]{\linewidth}\raggedright
\textbf{Class}
\end{minipage} & \begin{minipage}[b]{\linewidth}\raggedright
\textbf{Responsibility}
\end{minipage} \\
\midrule\noalign{}
\endhead
\bottomrule\noalign{}
\endlastfoot
Customer & One record per person, identified by a unique phone number;
owns bills, orders and at most one khata account. \\
KhataAccount / LedgerEntry & Credit limit and days; append-only entries
with running balance. Mistakes are fixed by a reversal entry, never by
editing. \\
Bill / BillItem / BillPayment & Shop sale with a unique bill number and
idempotency key; source tells POS bills from Tally imports; payment
split by mode. \\
Product / Variant & Catalog; variant holds price, MRP, GST, HSN, barcode
and the fall-pico flag. \\
StockLevel / StockMovement & Current on-hand and reserved quantity;
every change is recorded as a movement (one stock ledger). \\
Order / Payment & Online order lifecycle; payment confirmed only by
verified gateway webhook with a unique gateway ID. \\
FallPicoJob / FpWorker & Job state machine, worker assignment and
per-piece wages. \\
Employee / SalaryRun & Attendance-based salary, advances and
owner-approved monthly runs. \\
OutboxItem / MessageJob & Reliable delivery to Tally and WhatsApp: saved
in the same transaction as the sale and sent later with retries. \\
\end{longtable}

\subsection{Relationships Used}\label{relationships-used}

\begin{itemize}
\item
  \textbf{Composition} (filled diamond): Bill to BillItem and
  BillPayment, KhataAccount to LedgerEntry, Product to Variant; parts
  cannot exist without the whole.
\item
  \textbf{Aggregation} (hollow diamond): Customer to KhataAccount; the
  account belongs to a customer, who can exist without credit.
\item
  \textbf{Association} with multiplicity: Customer places 0..* Bills and
  Orders; a BillItem refers to one Variant and may create 0..1
  FallPicoJob; an FpWorker works on many jobs.
\end{itemize}

\subsection{Conclusion}\label{conclusion-3}

\begin{quote}
The class diagram turns the SRS nouns into a clean domain model. Key
design rules are visible directly in the model: money in paise as
integers, append-only ledgers, idempotency keys on bills and the outbox
for reliable integration.
\end{quote}

\clearpage
\section{LAB EXPERIMENT 5}\label{lab-experiment-5}

\subsection{Sequence Diagrams}\label{sequence-diagrams}

\begin{quote}
\textbf{Aim}

To show the time-ordered interaction between objects for the two most
important scenarios of Wereverse using sequence diagrams.
\end{quote}

\subsection{Sequence Diagram 1: Khata Sale at the POS to WhatsApp and
Tally (Mode
B)}\label{sequence-diagram-1-khata-sale-at-the-pos-to-whatsapp-and-tally-mode-b}

\includegraphics[width=6.22271in,height=6.14125in]{media/media/image5.jpeg}

\begin{quote}
\emph{Figure 5.1: Sequence diagram for a khata sale}

The POS saves the bill locally first and sends it with an idempotency
key, so a retry never creates a second bill. The API saves the bill,
stock change, ledger entry, fall-pico jobs and outbox rows in
\textbf{one database transaction}. Workers then deliver the WhatsApp
message and the Tally voucher asynchronously, so a WhatsApp or Tally
outage never blocks billing. The AI service only writes wording around
placeholders; the validator fills real numbers from the ledger.
\end{quote}

\subsection{Sequence Diagram 2: Customer Pays Dues through the Pay Now
Link}\label{sequence-diagram-2-customer-pays-dues-through-the-pay-now-link}

\includegraphics[width=6.0825in,height=5.23125in]{media/media/image6.jpeg}

\begin{quote}
\emph{Figure 5.2: Sequence diagram for Pay Now payment}

Payment is never trusted from the browser. Only the signed Razorpay
webhook creates the PAYMENT ledger entry, and a duplicate webhook with
the same payment ID does nothing. The payment is allocated to the oldest
bills first so the ageing report stays accurate.
\end{quote}

\subsection{Conclusion}\label{conclusion-4}

\begin{quote}
The sequence diagrams make the order of messages and the transaction
boundaries explicit. They show how the outbox pattern, idempotency keys
and webhook verification work together to keep money and stock correct.
\end{quote}

\clearpage
\section{LAB EXPERIMENT 6}\label{lab-experiment-6}

\subsection{Activity Diagram and State Chart
Diagram}\label{activity-diagram-and-state-chart-diagram}

\begin{quote}
\textbf{Aim}

To model the workflow of the khata billing process with an activity
diagram, and the life cycle of a fall-pico job with a state chart
diagram.
\end{quote}

\subsection{Activity Diagram: Khata Billing at the
Counter}\label{activity-diagram-khata-billing-at-the-counter}

\includegraphics[width=4.39969in,height=6.365in]{media/media/image7.png}

\begin{quote}
\emph{Figure 6.1: Activity diagram for khata billing}

Decision nodes handle repeated scanning, the credit-limit check with
manager approval, and the online/offline branch. The fork shows the
three actions done in the same transaction (bill and ledger, fall-pico
jobs, outbox rows), which join before the bill is printed and the
WhatsApp bill is sent.
\end{quote}

\subsection{State Chart Diagram: Fall-Pico
Job}\label{state-chart-diagram-fall-pico-job}

\includegraphics[width=6.32667in,height=1.015in]{media/media/image8.jpeg}

\begin{quote}
\emph{Figure 6.2: State chart diagram of a fall-pico job}
\end{quote}

\begin{longtable}[]{@{}
  >{\raggedright\arraybackslash}p{(\columnwidth - 2\tabcolsep) * \real{0.3000}}
  >{\raggedright\arraybackslash}p{(\columnwidth - 2\tabcolsep) * \real{0.7000}}@{}}
\toprule\noalign{}
\begin{minipage}[b]{\linewidth}\raggedright
\textbf{State}
\end{minipage} & \begin{minipage}[b]{\linewidth}\raggedright
\textbf{Meaning}
\end{minipage} \\
\midrule\noalign{}
\endhead
\bottomrule\noalign{}
\endlastfoot
PENDING & Created when a saree above the price limit is billed (or an
online order is paid) \\
ASSIGNED / IN\_PROGRESS & Supervisor assigns a worker; worker starts the
work \\
COMPLETED & Worker returns the saree for quality check \\
REWORK & QC failed; rework is not paid again \\
QC\_PASSED & Quality passed; wage = worker rate x pieces becomes
payable \\
READY / HANDED\_OVER & Customer told on WhatsApp; saree collected or
online order packed \\
CANCELLED & Customer declined or the bill was returned \\
\end{longtable}

\subsection{Conclusion}\label{conclusion-5}

\begin{quote}
The activity diagram captures the business rules of a khata sale,
including approvals and offline behaviour. The state chart gives the
fall-pico job a strict life cycle, which makes wages, tracking of every
saree and the online packing rule easy to enforce.
\end{quote}

\clearpage
\section{LAB EXPERIMENT 7}\label{lab-experiment-7}

\subsection{ER Diagram, Deployment Diagram and Test
Cases}\label{er-diagram-deployment-diagram-and-test-cases}

\begin{quote}
\textbf{Aim}

To design the data model (ER diagram), the physical deployment of
Wereverse, and the test cases used to verify the system.
\end{quote}

\subsection{ER Diagram}\label{er-diagram}

\includegraphics[width=5.82187in,height=5.56875in]{media/media/image9.png}

\begin{quote}
\emph{Figure 7.1: ER diagram of the main Wereverse entities}

Database rules: UUID primary keys; money stored in paise as integers;
dates in UTC shown in IST; financial and stock tables are append-only
with reversals only; unique constraints on bill number per counter,
voucher reference, gateway payment ID and idempotency keys; every bill
and stock movement records its source (pos, website, tally\_import).
\end{quote}

\subsection{Deployment Diagram}\label{deployment-diagram}

\includegraphics[width=6.19875in,height=3.47344in]{media/media/image10.png}

\begin{quote}
\emph{Figure 7.2: Deployment diagram of Wereverse}

The application runs on AWS Mumbai as a FastAPI modular monolith with
PostgreSQL, Redis and Celery workers, plus a separate AI service. The
POS is an installable web app on the counter PC. TallyPrime and the sync
agent run on a 24-hour cloud server in India; the agent only makes
\textbf{outbound} HTTPS calls, so Tally is never exposed to the
internet.
\end{quote}

\subsection{Test Cases}\label{test-cases}

\begin{longtable}[]{@{}
  >{\raggedright\arraybackslash}p{(\columnwidth - 6\tabcolsep) * \real{0.0700}}
  >{\raggedright\arraybackslash}p{(\columnwidth - 6\tabcolsep) * \real{0.1300}}
  >{\raggedright\arraybackslash}p{(\columnwidth - 6\tabcolsep) * \real{0.4000}}
  >{\raggedright\arraybackslash}p{(\columnwidth - 6\tabcolsep) * \real{0.4000}}@{}}
\toprule\noalign{}
\textbf{ID} & \textbf{Area} & \textbf{Test} & \textbf{Expected result} \\
\midrule\noalign{}
\endhead
\bottomrule\noalign{}
\endlastfoot
T01 & POS & 3-item cash bill & Bill printed, website stock drops, Tally
sales voucher within 5 minutes \\
T02 & POS & Split payment: cash + UPI + khata & Each part recorded; khata
part added to ledger \\
T04 & POS & Internet off, make 5 bills, turn on & All 5 synced once; no
duplicate numbers \\
T05 & Stock & Last unit bought at shop and online together & One
succeeds; the other gets a clear message \\
T06 & Khata & Credit sale within limit & Ledger entry; WhatsApp bill
with correct amounts within 1 minute \\
T07 & Khata & Credit sale over limit & Blocked; manager override logged
with reason \\
T08 & WhatsApp & AI text has a wrong amount & Standard template sent
with correct amount \\
T10 & Pay Now & Customer pays part amount & Balance falls; WhatsApp
receipt; Tally receipt voucher \\
T11 & Pay Now & Same webhook delivered twice & Only one ledger entry \\
T13 & Ledger & Try to edit or delete an entry & Denied; reversal works
and both entries show \\
T15 & Tally & Tally closed, 20 bills made, Tally opened & 20 vouchers
synced once each, in order \\
T19 & Payroll & 24 payable days, salary 15,000, advance 5,000 with
1,000 recovery, 6,000 paid & Earned 12,000; net 11,000; pending 5,000;
advance balance 4,000 \\
T22 & Fall-pico & Saree at 1,000 and at 1,001 & No job at 1,000; job at
1,001 \\
T24 & Fall-pico & Online saree order & Cannot be packed until QC passed
or skipped \\
T29 & Access & Cashier tries a write-off & Denied (403) \\
T34 & Mode A & Sale made in TallyPrime & Website stock drops within 5
minutes; sale saved once \\
T36 & Mode A & Khata bill with an unmatched party & Goes to Sync
Issues; no message sent \\
\end{longtable}

\begin{quote}
Release gate: no Critical or High bugs open; all test cases pass on
staging; a two-week parallel run shows zero unexplained difference
between Wereverse and Tally; owner, accountant and shop staff sign off
UAT.
\end{quote}

\subsection{Conclusion}\label{conclusion-6}

\begin{quote}
The ER diagram fixes the persistent structure behind the class model,
the deployment diagram shows where each component runs and how Tally
stays secure, and the test cases trace back to the requirements of
Experiment 1, with special focus on money, stock and
duplicate-prevention.
\end{quote}

\clearpage
\section{LAB EXPERIMENT 8}\label{lab-experiment-8}

\subsection{Modular Design of a Software
System}\label{modular-design-of-a-software-system}

\begin{quote}
\textbf{Aim}

To divide Wereverse into modules with one clear responsibility each,
define the interfaces between the modules, and justify the design in
terms of cohesion and coupling.
\end{quote}

\subsection{Background}\label{background}

\begin{quote}
\textbf{Modular design} divides a system into separate parts called
modules, each handling one specific responsibility, so that every module
can be built, tested and changed on its own. An \textbf{interface} is
the defined way in which one module asks another for a service: it
states the operation, the input and the output. \textbf{Cohesion}
measures how closely the tasks inside one module are related; high
cohesion is desirable. \textbf{Coupling} measures how strongly one
module depends on another; low coupling is desirable.

Wereverse is built as a \textbf{modular monolith} (SRS section 4.3): one
FastAPI application deployed as one unit and divided inside into
modules, each with its own router, schemas, service, repository, models
and events. Routers never write to tables directly; anything that writes
to stock, ledgers, salary or the outbox goes through a service function
with one transaction and an audit entry (SRS section 11.2). This keeps
the system simple for a five-member team, and modules can be split into
services later.
\end{quote}

\subsection{Functional Requirements}\label{functional-requirements-1}

\begin{quote}
Table 8.1 lists the functional requirements of Experiment 1 by their SRS
module codes. They are used below to check that every requirement
belongs to at least one module.

\emph{Table 8.1: Functional requirements of Wereverse}
\end{quote}

\begin{longtable}[]{@{}
  >{\raggedright\arraybackslash}p{(\columnwidth - 2\tabcolsep) * \real{0.1934}}
  >{\raggedright\arraybackslash}p{(\columnwidth - 2\tabcolsep) * \real{0.8066}}@{}}
\toprule\noalign{}
\textbf{ID} & \textbf{Requirement / action performed by the system} \\
\midrule\noalign{}
\endhead
\bottomrule\noalign{}
\endlastfoot
\textbf{FR-AUT, FR-ADM} & Role-based access, 2FA for the owner, complete
audit log (who, what, old value, new value, when, device) and
settings. \\
\textbf{FR-CAT} & Products and variants with barcodes and labels; one
stock ledger where every change is a stock movement; online reservation
released after 30 minutes; stock count with reasoned adjustments. \\
\textbf{FR-PUR} & Supplier bills, stock-in and label printing. \\
\textbf{FR-POS} & Scan-first billing; split payment (cash, UPI, card,
khata); discounts within role limits and manager PIN above; returns and
exchanges; cash register; offline billing with idempotent sync. \\
\textbf{FR-KHT} & Khata accounts with credit limit; append-only ledger;
credit-limit check with manager override; oldest-bill-first allocation;
ageing report; reminder rules. \\
\textbf{FR-WAP, FR-NOT} & WhatsApp bill, receipt, reminder and
fall-pico-ready messages from approved templates; AI wording around
placeholders only; SMS fallback; opt-in and quiet hours;
notifications. \\
\textbf{FR-PAY, FR-PRT} & Signed Pay Now link valid 7 days; Razorpay
payment; payment confirmed only by verified webhook; overpayment becomes
advance credit. \\
\textbf{FR-TLY} & Mode A import and Mode B push of sales; online orders
and receipts always pushed; item and ledger mapping; Sync Issues screen;
nightly reconciliation; unique voucher reference. \\
\textbf{FR-EMP} & Attendance, monthly and daily salary, incentive,
overtime, advances and recovery; salary run draft, owner approval and
lock; employee sees own statement only. \\
\textbf{FR-FPC} & Job for each eligible saree above the price limit;
tag; assignment; QC pass or fail with rework; per-piece wage on QC pass;
online order not packed until QC passed or skipped. \\
\textbf{FR-STO, FR-CHK, FR-ORD} & Browse, search, cart, checkout, order
tracking, shipping and returns; no khata option for online customers. \\
\textbf{FR-RPT} & Owner dashboard, sales and GST reports. \\
\textbf{FR-AI} & AI wording (Phase 1); recommendations, chatbot, smart
search, size help, descriptions, posters (Phase 2); visual search,
demand forecast, what-to-buy (Phase 3). \\
\end{longtable}

\subsection{Module Identification}\label{module-identification}

\begin{quote}
The modules follow the business domains of the shop. Table 8.2 gives the
main responsibility of each module and the requirements it covers.

\emph{Table 8.2: Modules of Wereverse and their responsibilities}
\end{quote}

\begin{longtable}[]{@{}
  >{\raggedright\arraybackslash}p{(\columnwidth - 4\tabcolsep) * \real{0.2210}}
  >{\raggedright\arraybackslash}p{(\columnwidth - 4\tabcolsep) * \real{0.5912}}
  >{\raggedright\arraybackslash}p{(\columnwidth - 4\tabcolsep) * \real{0.1878}}@{}}
\toprule\noalign{}
\textbf{Module} & \textbf{Main responsibility} & \textbf{Requirements covered} \\
\midrule\noalign{}
\endhead
\bottomrule\noalign{}
\endlastfoot
\textbf{M1 Access, Audit and Settings} & Login (password for staff, OTP
for customers), token refresh, roles and permission checks on every
route, 2FA for owner and admin, append-only audit log, shop settings
(fall-pico limit, tally\_mode, quiet hours). & FR-AUT, FR-ADM \\
\textbf{M2 Catalog, Stock and Purchasing} & Products, variants, barcodes
and labels; stock\_movements and stock\_levels; reservations for unpaid
online orders; stock count; suppliers, purchases and stock-in. & FR-CAT,
FR-PUR \\
\textbf{M3 POS Billing} & Scan-first bills, split payment, discounts and
manager PIN, returns, cash register, GST on the server, bill series per
counter, offline queue with idempotency key. & FR-POS \\
\textbf{M4 Customers and Khata Ledger} & Customers, khata accounts and
limits, append-only ledger (SALE, PAYMENT, RETURN, ADJUSTMENT,
WRITE\_OFF, REVERSAL), running balance, oldest-bill-first allocation,
ageing, reminder rules, daily integrity check. & FR-KHT \\
\textbf{M5 Online Store and Orders} & Browse, search, cart, checkout,
order status, courier shipments, returns and refunds, coupons. & FR-STO,
FR-CHK, FR-ORD \\
\textbf{M6 Payments and Pay Now} & Signed Pay Now links, Razorpay
orders, signature-verified webhooks, idempotency by gateway payment ID,
advance credit on overpayment. & FR-PAY, FR-PRT \\
\textbf{M7 WhatsApp Messaging} & Message jobs, approved templates,
validator for numbers, opt-in and opt-out, quiet hours, SMS fallback,
delivery status updates, notifications. & FR-WAP, FR-NOT \\
\textbf{M8 Tally Integration} & Outbox and inbox, mapping tables,
voucher types, Mode A import and Mode B push, Sync Issues, nightly
reconciliation, endpoints for the sync agent (API key). & FR-TLY \\
\textbf{M9 Employees and Payroll} & Employees, attendance, salary
formula, incentive, overtime, advances and recovery, salary run (draft,
approved, locked), own statement. & FR-EMP \\
\textbf{M10 Fall-Pico} & Job creation for eligible sarees, job states,
tags, assignment, QC and rework, worker wage ledger, monthly cost
report. & FR-FPC \\
\textbf{M11 Reports and Dashboard} & Owner dashboard, sales and GST
reports; read-only queries over other modules. & FR-RPT \\
\textbf{M12 AI Service} & Separate FastAPI service: message wording,
recommendations, chatbot, search, size help, descriptions, posters,
forecasts, each with a fallback. & FR-AI \\
\end{longtable}

\begin{quote}
\textbf{Mapping to the planned repository (SRS section 11.1):} M1 =
modules auth and settings; M2 = catalog, inventory and purchasing; M3 =
pos; M4 = customers and khata; M5 = orders and shipping, with the
storefront in apps/web; M6 = payments; M7 = notifications with the
whatsapp and sms integrations; M8 = tally with apps/tally-bridge; M9 =
employees; M10 = fall\_pico; M11 = analytics; M12 = apps/ai-service. The
Celery worker (apps/worker) runs outbox, reminder, Tally sync and
reconciliation jobs for M7 and M8.
\end{quote}

\subsection{Modular Design Diagram}\label{modular-design-diagram}

\begin{quote}
\begin{center}
\includegraphics[width=5.63542in,height=7.53819in]{media/media/image11.png}
\end{center}

Figure 8.1 shows the modules, rotated for readability. Solid arrows are
direct calls (service function or REST); dashed arrows are asynchronous
work that goes through the outbox and the worker.

\emph{Figure 8.1: Modular design of Wereverse (module numbers match
Table 8.2), rotated for readability}
\end{quote}

\subsection{Interfaces Between
Modules}\label{interfaces-between-modules}

\begin{quote}
Table 8.3 describes the main interfaces, with the input each accepts and
the output it returns. Money is always passed as whole paise, and errors
use the common codes of SRS section 8.2.

\emph{Table 8.3: Interfaces between modules of Wereverse}
\end{quote}

\begin{longtable}[]{@{}
  >{\raggedright\arraybackslash}p{(\columnwidth - 6\tabcolsep) * \real{0.1823}}
  >{\raggedright\arraybackslash}p{(\columnwidth - 6\tabcolsep) * \real{0.2597}}
  >{\raggedright\arraybackslash}p{(\columnwidth - 6\tabcolsep) * \real{0.2486}}
  >{\raggedright\arraybackslash}p{(\columnwidth - 6\tabcolsep) * \real{0.3094}}@{}}
\toprule\noalign{}
\textbf{Caller to provider} & \textbf{Operation / interface} & \textbf{Input} & \textbf{Output} \\
\midrule\noalign{}
\endhead
\bottomrule\noalign{}
\endlastfoot
M1 to every module & Permission check on every route & Bearer token from
the Authorization header & User id and role; or 401 UNAUTHORIZED (not
logged in or token expired) or 403 FORBIDDEN (role not allowed). \\
M3 to M2 & Deduct stock (service function, same transaction as the bill)
& variant\_id, qty, source = pos & Stock movement saved and on\_hand
reduced. An offline sale that makes stock negative is still accepted and
a Stock Conflict alert is raised. \\
M3 to M4 & Check credit limit and add SALE ledger entry & account\_id,
amount\_paise, bill id & New running balance; or 422
CREDIT\_LIMIT\_EXCEEDED unless a manager overrides. \\
M3 to M10 & Create fall-pico job for an eligible saree line & bill item
id, price, fall\_pico\_eligible flag & Job id (state PENDING) and a tag;
no job when the price is not above the limit or the customer
declines. \\
M3, M5, M6, M2, M9 to Outbox & Write outbox rows in the same transaction
& Event type, reference ids, unique voucher reference & Row with status
pending; no direct reply. The worker delivers it later with retries. \\
Outbox to M7 to WhatsApp provider & Send bill, receipt or reminder &
Purpose, facts (bill no., amount, balance), customer language & Message
job sent, delivered, read or failed; on failure an SMS is sent and the
ledger entry is flagged. \\
M7 to M12 & POST /ai/message/compose & purpose, language, tone,
first\_name, placeholders & Text containing only the placeholders and
source = ai; on timeout (3 s) or failed validation the standard approved
template is used. \\
M5 to M2 & Reserve and release stock for an online order & variant\_id,
qty & reserved increased; released on cancel or after 30 minutes unpaid;
409 OUT\_OF\_STOCK if the last unit is gone. \\
M5 to M6 & Start online payment & order id, amount\_paise & Razorpay
order id; the order stays PENDING\_PAYMENT until the verified webhook
arrives. \\
Razorpay to M6 & POST /webhooks/razorpay & Signature header, gateway
payment id, amount & Payment confirmed and order or ledger updated;
duplicate webhooks do nothing; a wrong signature is rejected. \\
M6 to M4 & Record PAYMENT entry and allocate oldest bill first &
account\_id, amount\_paise, gateway payment id & Allocations and new
balance; any overpayment becomes advance credit. \\
M5 to M10 & Pack check for a paid online order & order id & Packing
allowed only when every saree job is QC\_PASSED or skipped. \\
Tally agent to M8 & GET /tally/outbox/pending, POST
/tally/outbox/:id/result, POST /tally/inbox & Agent API key; success or
error result; raw vouchers read from Tally & Pending voucher items,
acknowledgement, and imported vouchers matched by the mapping tables. \\
M8 to M2 and M4 & Import sales, returns, purchases and stock from Tally
& Voucher data and mapping tables & Sale, ledger and stock rows with
source = tally\_import; an unmatched party or item goes to Sync Issues
and no WhatsApp message is sent. \\
M11 to M2 to M6 & Read-only report queries & Date range and filters &
Totals and lists; never writes. \\
\end{longtable}

\subsection{Cohesion Analysis}\label{cohesion-analysis}

\begin{quote}
Table 8.4 explains why the tasks inside the main modules belong
together.

\emph{Table 8.4: Cohesion of selected modules}
\end{quote}

\begin{longtable}[]{@{}
  >{\raggedright\arraybackslash}p{(\columnwidth - 4\tabcolsep) * \real{0.2099}}
  >{\raggedright\arraybackslash}p{(\columnwidth - 4\tabcolsep) * \real{0.3204}}
  >{\raggedright\arraybackslash}p{(\columnwidth - 4\tabcolsep) * \real{0.4696}}@{}}
\toprule\noalign{}
\textbf{Module} & \textbf{Tasks inside it} & \textbf{Why this has high cohesion} \\
\midrule\noalign{}
\endhead
\bottomrule\noalign{}
\endlastfoot
\textbf{M1 Access, Audit and Settings} & Login, token, roles, permission
check, 2FA, audit log, settings. & Every task answers one question: who
is this user and what may they do. It has no billing, stock or ledger
logic, so a change to login rules never touches other modules. \\
\textbf{M2 Catalog, Stock and Purchasing} & Products, variants, stock
movements, reservations, stock count, supplier purchases. & Every task
changes or describes stock. All stock changes are rows of one ledger, so
stock rules exist in exactly one place. \\
\textbf{M4 Customers and Khata Ledger} & Customers, khata accounts,
ledger entries, allocation, ageing, reminder rules. & All tasks work on
one thing: what a customer owes. Balance, allocation and reversal rules
are written once; POS and Payments only call this module. \\
\textbf{M7 WhatsApp Messaging} & Message jobs, templates, validator,
opt-in, quiet hours, SMS fallback. & Everything that talks to customers
by message is in one place. If the provider changes, only this module
changes. \\
\textbf{M8 Tally Integration} & Outbox and inbox, mapping, voucher
types, Sync Issues, reconciliation. & All knowledge of Tally and its
voucher format is kept here. Other modules only write outbox rows and
never build Tally XML. \\
\textbf{M9 Employees and Payroll} & Attendance, salary formula,
advances, recovery, salary run, statement. & All tasks serve one
purpose: pay staff correctly. The salary formula is defined once and the
run is locked after owner approval. \\
\textbf{M10 Fall-Pico} & Job creation, states, tags, QC, rework, worker
wages, cost. & One life-cycle (PENDING to HANDED\_OVER) and one wage
rule (payable once at QC\_PASSED). \\
\end{longtable}

\begin{quote}
Three main cohesion decisions were taken. First, all khata balance rules
are in M4, while M3 (POS) and M6 (Payments) only request an entry.
Second, M9 (monthly salary) and M10 (per-piece fall-pico wages) use the
same append-only ledger design but are separate modules, because their
rules differ (salary formula against wage after QC\_PASSED). Third,
wording by AI (M12) is kept apart from messaging (M7): the AI only
writes text around placeholders, and the numbers are filled and checked
by M7 from the database.
\end{quote}

\subsection{Coupling Analysis}\label{coupling-analysis}

\begin{quote}
Table 8.5 shows the main connections between modules, whether each is
needed, and how the coupling was kept low.

\emph{Table 8.5: Coupling between modules and techniques used to reduce
it}
\end{quote}

\begin{longtable}[]{@{}
  >{\raggedright\arraybackslash}p{(\columnwidth - 4\tabcolsep) * \real{0.1989}}
  >{\raggedright\arraybackslash}p{(\columnwidth - 4\tabcolsep) * \real{0.2431}}
  >{\raggedright\arraybackslash}p{(\columnwidth - 4\tabcolsep) * \real{0.5580}}@{}}
\toprule\noalign{}
\textbf{Module connection} & \textbf{Needed? Why?} & \textbf{How coupling is kept low} \\
\midrule\noalign{}
\endhead
\bottomrule\noalign{}
\endlastfoot
Clients to the API & Yes. One public API with token checking. & Clients
use only REST and JSON in one response format (success, data, error).
The POS keeps its own local copy and talks to the API only through
sync. \\
M3 to M2, M4, M10 & Yes. A bill must change stock and ledger and create
jobs together. & One transaction in a service function; only ids and
paise amounts are passed; M3 never reads their tables; the idempotency
key makes a retry safe. This is the tightest link and is kept on purpose
for correct money and stock. \\
M3, M5, M6, M9 to M7 and M8 & Only to notify or sync. & Outbox pattern:
the sender writes an outbox row and a worker delivers it with retries. A
WhatsApp or Tally outage never blocks a sale and no sale is lost. \\
M7 to M12 (AI) & Only for wording. & Separate service with 3 s timeout
and a standard-template fallback. Amounts always come from the database,
and no personal data is sent to the model. \\
M6 to Razorpay, M7 to WhatsApp provider & Yes. & Provider code sits in
app/integrations behind small client functions; a provider can be
replaced without changing the modules. Webhooks are signature-verified
and idempotent. \\
M8 to the Tally agent & Yes. & The agent makes outbound HTTPS calls
only, with its own API key. Tally is never opened to the internet. Each
voucher has a unique reference, so a retry cannot duplicate it. \\
M9 and M10 & No direct link. & Separate tables and rules. They share
only the design of an append-only ledger, not code or data. \\
M5 to M10 & Weakly (pack check). & M5 only asks for the state of the
jobs of one order. \\
M11 to other modules & Read only. & Reports never write, so they can be
changed or removed without affecting billing. \\
\end{longtable}

\begin{quote}
General techniques used: modules exchange ids, paise amounts and small
payloads, not internal objects; money, stock and ledger tables are
append-only; shared packages hold only types and configuration.

\textbf{Known remaining coupling.} All modules share one PostgreSQL
database, and some tables refer to others by id (for example bills to
customers). M11 reads tables of several modules. M3, M2, M4 and M10 must
be deployed together because they share one transaction. This is
accepted for the modular monolith chosen in the SRS; a later improvement
is to let M11 read through database views and to split modules into
services only if the shop grows.
\end{quote}

\subsection{Design Review Checklist}\label{design-review-checklist}

\begin{quote}
\emph{Table 8.6: Design review checklist for Experiment 8}
\end{quote}

\begin{longtable}[]{@{}
  >{\raggedright\arraybackslash}p{(\columnwidth - 2\tabcolsep) * \real{0.6409}}
  >{\raggedright\arraybackslash}p{(\columnwidth - 2\tabcolsep) * \real{0.3591}}@{}}
\toprule\noalign{}
\textbf{Review item} & \textbf{Status} \\
\midrule\noalign{}
\endhead
\bottomrule\noalign{}
\endlastfoot
Every important requirement belongs to at least one module. & Satisfied
(Table 8.2) \\
Every module has a clear purpose. & Satisfied (Tables 8.2 and 8.4) \\
Diagram matches the module table. & Satisfied (Figure 8.1) \\
Interfaces show inputs and outputs. & Satisfied (Table 8.3) \\
Unnecessary direct dependencies are avoided. & Satisfied (Table 8.5) \\
Remaining coupling is stated honestly. & Satisfied (note after Table
8.5) \\
\end{longtable}

\subsection{Results and Discussion}\label{results-and-discussion}

\begin{quote}
Wereverse was divided into eleven business modules (M1 to M11) plus a
separate AI service (M12). Every functional requirement is covered by at
least one module (Table 8.2), and the diagram agrees with the table
(Figure 8.1). Cohesion is high because each business domain lives in one
module with one purpose: all khata rules in M4, all Tally knowledge in
M8, all salary rules in M9. Coupling is low because modules exchange
only ids and small payloads, slow or unreliable work (WhatsApp, Tally,
AI) is decoupled through the outbox, timeouts and fallbacks, and every
provider sits behind its own integration code. The review also shows the
coupling that remains by design: the shared database and the single
transaction of a bill.
\end{quote}

\subsection{Conclusion}\label{conclusion-7}

\begin{quote}
Wereverse was divided into twelve modules with defined interfaces, and
the design was justified through cohesion and coupling. A good modular
design places each business domain in one module and lets modules
communicate through small, well-defined interfaces. The few strong
links, such as the bill transaction, are kept on purpose where
correctness of money and stock matters.
\end{quote}

\clearpage
\section{LAB EXPERIMENT 9}\label{lab-experiment-9}

\subsection{Functional Test Case
Design}\label{functional-test-case-design}

\begin{quote}
\textbf{Aim}

To design functional test cases for selected Wereverse requirements,
write clear expected results, and link each test case to its
requirement.
\end{quote}

\subsection{Background}\label{background-1}

\begin{quote}
\textbf{Functional testing} checks whether the software behaves as its
requirements state, by giving inputs and comparing the observed output
with the expected output. A \textbf{test case} has a unique ID, the
requirement it checks, the input or precondition, the steps, the
expected result and, after execution, the actual result and status. A
\textbf{positive test} uses valid input and expects normal behaviour; a
\textbf{negative test} uses invalid input and expects correct rejection.
\textbf{Boundary value testing} checks values at or just beyond a limit,
where errors are most common. A \textbf{traceability matrix} links every
requirement to its tests, so that no requirement is left untested.

\textbf{Stage of this experiment.} Wereverse is at the design stage and
has no running build yet. The test cases and their expected results are
therefore designed from the SRS (module notes in section 5, error codes
in section 8.2 and the worked examples of the test plan in section 13).
The Actual result and Status columns are prepared for the day the tests
are run; no execution result is claimed in this report.
\end{quote}

\subsection{Requirements Selected for
Testing}\label{requirements-selected-for-testing}

\begin{quote}
\emph{Table 9.1: Requirements selected for testing}
\end{quote}

\begin{longtable}[]{@{}
  >{\raggedright\arraybackslash}p{(\columnwidth - 4\tabcolsep) * \real{0.0994}}
  >{\raggedright\arraybackslash}p{(\columnwidth - 4\tabcolsep) * \real{0.6188}}
  >{\raggedright\arraybackslash}p{(\columnwidth - 4\tabcolsep) * \real{0.2818}}@{}}
\toprule\noalign{}
\textbf{Req. ID} & \textbf{Functional requirement in simple words} & \textbf{SRS reference} \\
\midrule\noalign{}
\endhead
\bottomrule\noalign{}
\endlastfoot
\textbf{R-01} & A user can log in with valid credentials; wrong or
missing data is rejected; protected APIs need a valid token and the
right role. & FR-AUT; SRS 8.2, 12.1; test T29 \\
\textbf{R-02} & A khata sale is allowed within the
customer\textquotesingle s credit limit, blocked above it, and can be
approved by a manager; a repeated bill is not created twice. & FR-KHT,
FR-POS; SRS 5.3, 8.4; tests T02, T06, T07 \\
\textbf{R-03} & Salary is calculated by the SRS formula; advances above
the limit need owner approval; salary data is visible only to allowed
roles; an approved run is locked. & FR-EMP; SRS 5.7; tests T19, T20,
T21 \\
\textbf{R-04} & A fall-pico job is created only for an eligible saree
priced above the limit; rework is not paid twice; an online saree is not
packed before QC. & FR-FPC; SRS 5.8; tests T22, T23, T24 \\
\end{longtable}

\subsection{Test Data and Assumptions}\label{test-data-and-assumptions}

\begin{quote}
The rules below are taken from the SRS. Test accounts and amounts are
test data chosen by the group; anything the SRS leaves open is marked as
an assumption.
\end{quote}

\begin{itemize}
\item
  \textbf{Environment:} staging system with test data; API base /api/v1;
  protected calls carry Authorization: Bearer token; money is stored in
  paise and shown here in rupees; errors use the codes of SRS section
  8.2.
\item
  \textbf{Accounts:} one test cashier, test manager, test owner and test
  employee.
\item
  \textbf{Khata test customer:} credit limit Rs. 10,000 and current
  balance Rs. 8,000.
\item
  \textbf{Credit limit boundary:} the POS warns when the limit is
  crossed. A bill that takes the balance exactly to the limit is treated
  as allowed {[}Assumption: the limit is inclusive; to be confirmed with
  the client{]}.
\item
  \textbf{Payroll:} payable days = present + half-day x 0.5 + paid leave
  + weekly off; monthly staff earned = salary x payable days / days in
  month; gross = earned + incentive + overtime; net = gross - advance
  recovery - other deductions; pending = net - salary already paid (SRS
  5.7).
\item
  \textbf{Fall-pico:} a job is created for each saree with
  fall\_pico\_eligible = true and a price above the limit (setting,
  default Rs. 1,000). Whether the limit applies before or after discount
  is an open decision in the SRS, so the tests use bills without
  discount. Worker rate used in the test: Rs. 60 per piece.
\end{itemize}

\subsection{Test Case Register}\label{test-case-register}

\begin{quote}
Table 9.2 contains 20 test cases with positive and negative cases and
boundary values. The expected results were written from the SRS rules.

\emph{Table 9.2: Test case register of Wereverse}
\end{quote}

\begin{longtable}[]{@{}
  >{\raggedright\arraybackslash}p{(\columnwidth - 6\tabcolsep) * \real{0.1050}}
  >{\raggedright\arraybackslash}p{(\columnwidth - 6\tabcolsep) * \real{0.2099}}
  >{\raggedright\arraybackslash}p{(\columnwidth - 6\tabcolsep) * \real{0.3149}}
  >{\raggedright\arraybackslash}p{(\columnwidth - 6\tabcolsep) * \real{0.3702}}@{}}
\toprule\noalign{}
\textbf{Test ID / Req.} & \textbf{Scenario and type} & \textbf{Test input / precondition} & \textbf{Expected result} \\
\midrule\noalign{}
\endhead
\bottomrule\noalign{}
\endlastfoot
TC-01 R-01 & Valid staff login (Positive) & POST /auth/login with the
test cashier\textquotesingle s correct credentials. & HTTP 200, success =
true; token returned; user role = Cashier; error = null. \\
TC-02 R-01 & Wrong password (Negative) & Test cashier\textquotesingle s
username with a wrong password. & HTTP 401 UNAUTHORIZED; no token
issued. \\
TC-03 R-01 & Missing password (Negative) & Username only; password
field not sent. & HTTP 400 VALIDATION\_ERROR; no token issued. \\
TC-04 R-01 & Protected API without token (Negative) & GET
/khata/accounts/:id/ledger with no Authorization header. & HTTP 401
UNAUTHORIZED; no ledger data returned. \\
TC-05 R-01 & Cashier tries a write-off (Negative, role check) &
Logged-in cashier calls POST /khata/accounts/:id/adjustments (owner
only). & HTTP 403 FORBIDDEN; no ledger entry is created. \\
TC-06 R-02 & Khata sale within limit (Positive) & Customer limit Rs.
10,000, balance Rs. 8,000. POST /pos/bills with a bill of Rs. 1,500
(150000 paise) fully on khata. & Bill saved; SALE ledger entry of 150000
paise; new balance Rs. 9,500 (950000 paise); WhatsApp job and Tally
outbox row queued. \\
TC-07 R-02 & Sale exactly at the limit (Boundary) & Same customer; bill
of Rs. 2,000 on khata, so the balance becomes exactly Rs. 10,000. & Bill
accepted; new balance Rs. 10,000 {[}Assumption: limit is
inclusive{]}. \\
TC-08 R-02 & Sale crosses the limit (Negative) & Same customer; bill of
Rs. 3,200 on khata, so the balance would be Rs. 11,200. & HTTP 422
CREDIT\_LIMIT\_EXCEEDED; message says the limit is crossed by Rs. 1,200
and manager approval is needed; no bill, ledger entry or stock
deduction. \\
TC-09 R-02 & Manager override (Positive) & Same bill as TC-08; manager
enters the PIN and a reason. & Bill saved; balance Rs. 11,200; override
logged with approver and reason. \\
TC-10 R-02 & Same bill sent twice (Negative, idempotency) & POST
/pos/bills twice with the same idempotency\_key (for example
C1-26-000123). & Second call returns the same bill\_id and response;
only one bill and one ledger entry exist; stock is deducted once. \\
TC-11 R-03 & Monthly salary calculation (Positive) & Salary Rs. 15,000,
30-day month; attendance: 20 present, 2 half-day, 1 paid leave, 2 weekly
off; advance Rs. 5,000 with recovery Rs. 1,000; Rs. 6,000 already
paid. & Payable days 24; earned Rs. 12,000; gross Rs. 12,000; net Rs.
11,000; pending Rs. 5,000; advance balance falls from Rs. 5,000 to Rs.
4,000. \\
TC-12 R-03 & Daily-rate staff with incentive and overtime (Positive) &
Daily rate Rs. 500; 22 payable days; incentive Rs. 500; overtime Rs.
300; no advance. & Earned Rs. 11,000; gross Rs. 11,800; net payable Rs.
11,800. \\
TC-13 R-03 & Advance above the limit (Negative) & An advance larger than
the configured advance limit is requested for an employee. & Advance is
not recorded until the owner approves it. \\
TC-14 R-03 & Salary access by role (Negative) & Test manager calls the
salary APIs of an employee; the test employee calls GET /me/statement. &
Manager: HTTP 403 FORBIDDEN. Employee: sees only their own
statement. \\
TC-15 R-03 & Change after approval (Negative) & Owner approves the
salary run; a salary line is then edited. & Change rejected with a
business-rule error (HTTP 422); the run stays locked. \\
TC-16 R-04 & Saree above the limit (Positive, boundary) & Eligible
saree priced Rs. 1,001, no discount; customer does not decline. & One
fall-pico job is created in state PENDING and a tag is printed. \\
TC-17 R-04 & Saree exactly at the limit (Boundary) & Eligible saree
priced Rs. 1,000, no discount. & No fall-pico job is created. \\
TC-18 R-04 & Saree not marked eligible (Negative) & Saree priced Rs.
1,500 with fall\_pico\_eligible = false. & No fall-pico job is
created. \\
TC-19 R-04 & QC fail, rework, then pass (Positive) & Job fails QC, goes
to REWORK, then passes QC; worker rate Rs. 60; 1 piece. & No wage after
the first failure; wage of Rs. 60 added once at QC\_PASSED; worker
statement total is Rs. 60. \\
TC-20 R-04 & Pack online saree before QC (Negative) & Paid online order
whose saree job is IN\_PROGRESS; staff tries to pack it, then again
after QC\_PASSED. & First attempt rejected; after QC\_PASSED (or
skipped) packing is allowed. \\
\end{longtable}

\subsection{Requirement Traceability
Check}\label{requirement-traceability-check}

\begin{quote}
Table 9.3 links each requirement to its test cases and confirms that
both positive and negative cases exist. Every requirement has at least
one positive and one negative test, and every test maps to a
requirement.

\emph{Table 9.3: Requirement traceability matrix}
\end{quote}

\begin{longtable}[]{@{}
  >{\raggedright\arraybackslash}p{(\columnwidth - 6\tabcolsep) * \real{0.1878}}
  >{\raggedright\arraybackslash}p{(\columnwidth - 6\tabcolsep) * \real{0.4641}}
  >{\raggedright\arraybackslash}p{(\columnwidth - 6\tabcolsep) * \real{0.1657}}
  >{\raggedright\arraybackslash}p{(\columnwidth - 6\tabcolsep) * \real{0.1823}}@{}}
\toprule\noalign{}
\textbf{Requirement ID} & \textbf{Related test case IDs} & \textbf{Positive case?} & \textbf{Negative case?} \\
\midrule\noalign{}
\endhead
\bottomrule\noalign{}
\endlastfoot
R-01 & TC-01, TC-02, TC-03, TC-04, TC-05 & Yes & Yes \\
R-02 & TC-06, TC-07, TC-08, TC-09, TC-10 & Yes & Yes \\
R-03 & TC-11, TC-12, TC-13, TC-14, TC-15 & Yes & Yes \\
R-04 & TC-16, TC-17, TC-18, TC-19, TC-20 & Yes & Yes \\
\end{longtable}

\subsection{Detailed Test Case
Records}\label{detailed-test-case-records}

\begin{quote}
Two test cases are recorded in full: TC-08 (a negative case) and TC-11
(a positive calculation case). Their actual result and status stay open
until the tests are run.

\emph{Table 9.4: Detailed record of test case TC-08}
\end{quote}

\begin{longtable}[]{@{}
  >{\raggedright\arraybackslash}p{(\columnwidth - 2\tabcolsep) * \real{0.2541}}
  >{\raggedright\arraybackslash}p{(\columnwidth - 2\tabcolsep) * \real{0.7459}}@{}}
\toprule\noalign{}
\textbf{Field} & \textbf{Details} \\
\midrule\noalign{}
\endhead
\bottomrule\noalign{}
\endlastfoot
\textbf{Test case ID / requirement ID} & TC-08 / R-02 \\
\textbf{Title and type} & Khata sale that crosses the credit limit -
Negative \\
\textbf{Preconditions} & Staging system is running; the test customer
has a khata account with limit Rs. 10,000 and balance Rs. 8,000; the
test cashier exists. \\
\textbf{Test input / test data} & Bill of Rs. 3,200 (320000 paise) paid
fully by khata; new idempotency\_key. \\
\textbf{Steps to follow} & 1. Log in as the test cashier with POST
/auth/login and copy the token. 2. Check the balance and limit with GET
/khata/accounts/:id/ledger. 3. Send POST /pos/bills with the khata
payment of 320000 paise and the Bearer token. 4. Read the HTTP status,
error code and message. 5. Check that no bill, ledger entry or stock
movement was created. \\
\textbf{Expected result} & HTTP 422 with code CREDIT\_LIMIT\_EXCEEDED
and a message that the limit is crossed by Rs. 1,200 and manager
approval is needed; balance stays Rs. 8,000; stock unchanged. \\
\textbf{Actual result} & To be recorded when the test is executed. \\
\textbf{Execution status} & Pending \\
\end{longtable}

\begin{quote}
\emph{Table 9.5: Detailed record of test case TC-11}
\end{quote}

\begin{longtable}[]{@{}
  >{\raggedright\arraybackslash}p{(\columnwidth - 2\tabcolsep) * \real{0.2541}}
  >{\raggedright\arraybackslash}p{(\columnwidth - 2\tabcolsep) * \real{0.7459}}@{}}
\toprule\noalign{}
\textbf{Field} & \textbf{Details} \\
\midrule\noalign{}
\endhead
\bottomrule\noalign{}
\endlastfoot
\textbf{Test case ID / requirement ID} & TC-11 / R-03 \\
\textbf{Title and type} & Monthly salary calculation with half-day, paid
leave and advance recovery - Positive \\
\textbf{Preconditions} & Staging system is running; the test employee
has a monthly salary of Rs. 15,000 and an advance of Rs. 5,000; the
month has 30 days. \\
\textbf{Test input / test data} & Attendance: 20 present, 2 half-day, 1
paid leave, 2 weekly off. Advance recovery Rs. 1,000. Salary already
paid Rs. 6,000. No incentive, overtime or other deductions. \\
\textbf{Steps to follow} & 1. As the test manager (with PIN), enter the
attendance with POST /employees/:id/attendance. 2. As the test owner,
record the advance recovery and the salary payment of Rs. 6,000 with
POST /employees/:id/salary-payments. 3. Create the salary run with POST
/payroll/runs. 4. Read payable days, earned, gross, net payable, pending
and advance balance. \\
\textbf{Expected result} & Payable days 24 (20 + 2 x 0.5 + 1 + 2);
earned Rs. 12,000 (15,000 x 24 / 30); gross Rs. 12,000; net payable Rs.
11,000; pending Rs. 5,000; advance balance Rs. 4,000. \\
\textbf{Actual result} & To be recorded when the test is executed. \\
\textbf{Execution status} & Pending \\
\end{longtable}

\subsection{Execution Log}\label{execution-log}

\begin{quote}
Table 9.6 is the log that will be filled in when the tests are run on
the staging system.

\emph{Table 9.6: Execution log of all test cases}
\end{quote}

\begin{longtable}[]{@{}
  >{\raggedright\arraybackslash}p{(\columnwidth - 6\tabcolsep) * \real{0.1105}}
  >{\raggedright\arraybackslash}p{(\columnwidth - 6\tabcolsep) * \real{0.4751}}
  >{\raggedright\arraybackslash}p{(\columnwidth - 6\tabcolsep) * \real{0.1657}}
  >{\raggedright\arraybackslash}p{(\columnwidth - 6\tabcolsep) * \real{0.2486}}@{}}
\toprule\noalign{}
\textbf{Test ID} & \textbf{Actual result observed} & \textbf{Status} & \textbf{Date and tester} \\
\midrule\noalign{}
\endhead
\bottomrule\noalign{}
\endlastfoot
TC-01 & & Pending & \\
TC-02 & & Pending & \\
TC-03 & & Pending & \\
TC-04 & & Pending & \\
TC-05 & & Pending & \\
TC-06 & & Pending & \\
TC-07 & & Pending & \\
TC-08 & & Pending & \\
TC-09 & & Pending & \\
TC-10 & & Pending & \\
TC-11 & & Pending & \\
TC-12 & & Pending & \\
TC-13 & & Pending & \\
TC-14 & & Pending & \\
TC-15 & & Pending & \\
TC-16 & & Pending & \\
TC-17 & & Pending & \\
TC-18 & & Pending & \\
TC-19 & & Pending & \\
TC-20 & & Pending & \\
\end{longtable}

\subsection{Review Checklist}\label{review-checklist}

\begin{quote}
\emph{Table 9.7: Review checklist for Experiment 9}
\end{quote}

\begin{longtable}[]{@{}
  >{\raggedright\arraybackslash}p{(\columnwidth - 2\tabcolsep) * \real{0.5967}}
  >{\raggedright\arraybackslash}p{(\columnwidth - 2\tabcolsep) * \real{0.4033}}@{}}
\toprule\noalign{}
\textbf{Review item} & \textbf{Status} \\
\midrule\noalign{}
\endhead
\bottomrule\noalign{}
\endlastfoot
Every test has a unique ID. & Satisfied (Table 9.2) \\
Every test maps to a requirement. & Satisfied (Table 9.3) \\
Both valid and invalid cases are included. & Satisfied (Table 9.2) \\
Expected results are clear. & Satisfied (Tables 9.2, 9.4 and 9.5) \\
Steps can be repeated by another person. & Satisfied for the two
detailed records (Tables 9.4 and 9.5) \\
All rules match the SRS; open points are marked. & Satisfied (Test Data
and Assumptions) \\
Tests have been executed. & Not yet: build pending (Table 9.6 is
ready) \\
\end{longtable}

\subsection{Results and Discussion}\label{results-and-discussion-1}

\begin{quote}
Twenty functional test cases were designed for four Wereverse
requirements: login and access control, khata credit limit, salary
calculation and fall-pico jobs. Each test is linked to a requirement and
every requirement has positive and negative tests (Table 9.3). Boundary
values were chosen where the SRS gives a limit: a sale that ends exactly
at the credit limit and one that crosses it, a saree priced Rs. 1,000
against Rs. 1,001, and 24 payable days from mixed attendance. The worked
salary example of the SRS (net Rs. 11,000, pending Rs. 5,000, advance
Rs. 4,000) is used as the expected result of TC-11.

Three points need a decision from the client: whether the credit limit
is inclusive (TC-07), whether the fall-pico limit applies before or
after discount (TC-16, TC-17), and the value of the advance limit
(TC-13). They are marked as assumptions and do not change the structure
of the tests.
\end{quote}

\clearpage
\section{OVERALL CONCLUSION}\label{overall-conclusion}

\begin{quote}
Through Experiments 1 to 9 we carried Wereverse from a business problem
to a complete design and a test plan. Requirements analysis (Experiment
1) set out what the shop needs and which qualities are non-negotiable:
billing that never stops, exact money, one shared stock and Tally books
that always match. Choosing an incremental process model (Experiment 2)
let us plan a rollout that keeps the live shop running while each module
is introduced and tested.

The UML diagrams each answered a different question. The use case
diagram showed who uses the system and for what; the class diagram fixed
the domain model and its rules; sequence diagrams exposed transaction
boundaries and the outbox pattern; activity and state chart diagrams
captured workflows and life cycles; and the ER and deployment diagrams
turned the design into storage and infrastructure. Modular design
(Experiment 8) divided the system into cohesive modules with small
interfaces and low coupling, and the functional test cases
(Experiment 9) closed the loop back to the requirements.

We learned that careful requirements and design work catches the hardest
problems (duplicate vouchers, offline conflicts, wrong amounts in
messages) before any code is written, which is cheaper and safer than
fixing them in production.
\end{quote}

\clearpage
\section{REFERENCES}\label{references}

\begin{enumerate}
\def\labelenumi{\arabic{enumi}.}
\item
  Wereverse Software Requirements Specification (SRS) v1.0 and Complete
  Project Documentation v1.0 (Draft), 5 October 2026.
\item
  I. Sommerville, \emph{Software Engineering}, 10th ed., Pearson, 2016.
\item
  R. S. Pressman and B. R. Maxim, \emph{Software Engineering: A
  Practitioner\textquotesingle s Approach}, 9th ed., McGraw-Hill, 2020.
\item
  IEEE Std 830-1998, \emph{IEEE Recommended Practice for Software
  Requirements Specifications}, IEEE, 1998.
\item
  Object Management Group, \emph{Unified Modeling Language (UML)
  Specification}, Version 2.5.1, 2017.
\item
  FastAPI documentation,
  \href{https://fastapi.tiangolo.com/}{https://fastapi.tiangolo.com}
\item
  PostgreSQL documentation, \url{https://www.postgresql.org/docs/}
\item
  Tally Solutions, TallyPrime integration (XML/HTTP) documentation,
  \href{https://help.tallysolutions.com/}{https://help.tallysolutions.com}
\item
  Razorpay payment gateway and webhooks documentation,
  \url{https://razorpay.com/docs/}
\item
  Meta, WhatsApp Business Platform documentation,
  \url{https://developers.facebook.com/docs/whatsapp}
\end{enumerate}

\end{document}